\documentclass[../../main.tex]{subfiles}

\begin{document}

\section{紧性与弱序列紧性: Eberlein-\v{S}mulian 定理}\label{section:ch15-eberlein-smulian}

\begin{theorem}[Goldstine]\label{theorem:ch15-goldstine}
令 $X$ 为赋范线性空间, $B$ 为 $X$ 的闭单位球, 而 $B^{**}$ 是 $X^{**}$ 的闭单位球. 则 $J(B)$ 的弱 $^{*}$ 闭包是 $B^{**}$.
\end{theorem}

\begin{proof}
根据前一章的推论 9, $J$ 是等距. 因此 $J(B)\subseteq B^{**}$. 令 $C$ 为 $J(B)$ 的弱 $^{*}$ 闭包. 我们把证明 $B^{**}$ 是弱 $^{*}$ 闭的留作练习. 因此 $C\subseteq B^{**}$. 由于 $B$ 是凸的而 $J$ 是线性的, $J(B)$ 是凸的. 前一章的命题 22 告诉我们, 在局部凸拓扑向量空间, 凸集的闭包是凸的. 因此 $C$ 关于弱 $^{*}$ 拓扑是闭的凸集. 假定 $C\neq B^{**}$. 令 $T$ 属于 $B^{**}\sim C$. 现在 $X$ 被 $(X^{*})^{*}$ 代替, 而考虑 $(X^{*})^{*}$ 为具有弱 $^{*}$ 拓扑的局部凸拓扑向量空间的情形下, 我们援引超平面分离定理, 见前一章的推论 26. 因此存在某个 $\psi\in X^{*}$ 使得 $\|\psi\|=1$ 且
\begin{align}
T(\psi)<\inf_{S\in C}S(\psi).\label{align:ch15-goldstine-sep}
\end{align}
观察到由于 $C$ 包含 $J(B)$,
\begin{align*}
\inf_{S\in C}S(\psi)\leqslant\inf_{x\in B}\psi(x)=-1.
\end{align*}
因此, $T(\psi)<-1$. 这是一个矛盾, 由于 $\|T\|\leqslant 1$ 而 $\|\psi\|=1$. 因此 $C=J(B)$, 证明完毕.
\end{proof}

\begin{lemma}\label{lemma:ch15-finite-dim-approx}
令 $X$ 为赋范线性空间而 $W$ 是 $X^{*}$ 的有限维子空间. 则存在 $X$ 的有限子集 $F$ 使得
\begin{align}
\text{对所有 }\psi\in W,\quad\|\psi\|/2\leqslant\max_{x\in F}\psi(x).\label{align:ch15-finite-dim}
\end{align}
\end{lemma}

\begin{proof}
由于 $W$ 是有限维的, 它的闭单位球 $S^{*}=\{\psi\in W\mid\|\psi\|=1\}$ 是紧的, 因此是全有界的. 选取 $S^{*}$ 的有限子集 $\{\psi_{1},\cdots,\psi_{n}\}$ 使得 $S^{*}\subseteq\bigcup_{k=1}^{n}B(\psi_{k},1/4)$. 对 $1\leqslant k\leqslant n$, 选取 $X$ 中的单位向量 $x_{k}$ 使得 $\psi_{k}(x_{k})>3/4$. 令 $\psi$ 属于 $S^{*}$. 观察到
\begin{align*}
\text{对 }1\leqslant k\leqslant n,\psi(x_{k})=\psi_{k}(x_{k})+[\psi-\psi_{k}](x_{k})\geqslant 3/4+[\psi-\psi_{k}](x_{k}).
\end{align*}
若我们选取 $k$ 使得 $\|\psi-\psi_{k}\|<1/4$, 则由于 $\|x_{k}\|=1$, $\psi(x_{k})\geqslant 1/2=1/2\|\psi\|$. 因此, 若 $F=\{x_{1},\cdots,x_{n}\}$ 且对 $\psi\in W$ 有 $\|\psi\|=1$, 则\eqref{align:ch15-finite-dim}成立. 因此它对所有 $\psi\in W$ 成立.
\end{proof}

\begin{theorem}[Eberlein-\v{S}mulian]\label{theorem:ch15-eberlein-smulian}
令 $B$ 为 Banach 空间 $X$ 的闭单位球. 则 $B$ 是弱紧的当且仅当它是弱序列紧的.
\end{theorem}

\begin{proof}
我们首先假设 $B$ 是紧的. Kakutani 定理告诉我们 $X$ 是自反的. 根据前一章的定理 17, $X$ 中的每个有界序列有弱收敛的子序列. 由于 $B$ 是弱闭的, $B$ 是弱序列紧的.

为证明反命题, 假设 $B$ 是弱序列紧的. 根据 Kakutani 定理, 为证明 $B$ 是紧的, 仅须证明 $X$ 是自反的. 令 $T$ 属于 $B^{**}$. Goldstine 定理告诉我们 $T$ 属于 $J(B)$ 的弱 $^{*}$ 闭包. 我们将用前一个引理证明 $T$ 属于 $J(B)$.

选取 $\psi_{1}\in B^{*}$. 由于 $T$ 属于 $J(B)$ 的弱 $^{*}$ 闭包, 可以选取 $x_{1}\in B$ 使得 $J(x_{1})$ 属于 $\mathcal{N}_{1,\psi_{1}}$. 定义 $N(1)=1$ 与 $W_{1}=\operatorname{span}[\{T,\ J(x_{1})\}]\subseteq X^{**}$. 令 $n$ 为自然数使得自然数 $N(n)$、$B$ 的子集 $\{x_{k}\}_{1\leqslant k\leqslant n}$、$X^{*}$ 的子集 $\{\psi_{k}\}_{1\leqslant k\leqslant N(n)}$ 已被定义, 且已定义 $W_{n}=\operatorname{span}[\{T,\ J(x_{1}),\cdots,J(x_{n})\}]$.

由于 $T$ 属于 $J(B)$ 的弱 $^{*}$ 闭包, 我们可以选取 $x_{(n+1)}\in B$ 使得
\begin{align}
J(x_{n+1})\in\mathcal{N}_{1/(n+1),\psi_{1},\cdots,\psi_{N(n)}}(T).\label{align:ch15-smulian-approx}
\end{align}
定义
\begin{align}
W_{n+1}=\operatorname{span}[\{T,\ J(x_{1}),\cdots,J(x_{n+1})\}].\label{align:ch15-smulian-W}
\end{align}
我们从前一个引理(其中 $X$ 被 $X^{*}$ 代替)推出, 存在自然数 $N(n+1)>N(n)$ 与 $X^{*}$ 的有限子集 $\{\psi_{k}\}_{N(n)<k\leqslant N(n+1)}$ 使得
\begin{align}
\text{对所有 }S\in W_{n+1},\quad\|S\|/2<\max_{N(n)<k\leqslant N(n+1)}S(\psi_{k}).\label{align:ch15-smulian-sep}
\end{align}
我们已归纳地定义了严格递增的自然数列 $\{N(n)\}$、$B$ 中的序列 $\{x_{n}\}$、$X^{*}$ 中的序列 $\{\psi_{n}\}$ 以及 $X^{**}$ 的子空间序列 $\{W_{n}\}$, 使得\eqref{align:ch15-smulian-approx}与\eqref{align:ch15-smulian-sep}成立. 由于 $\{W_{n}\}$ 是使得对每个指标 $n$\eqref{align:ch15-smulian-sep}成立的上升的序列,
\begin{align}
\text{对所有 }S\in W=\overline{\operatorname{span}}[\{T,\ J(x_{1}),\cdots,J(x_{n}),\cdots\}],\quad\|S\|/2\leqslant\sup_{1\leqslant k\leqslant\infty}S(\psi_{k}).\label{align:ch15-smulian-Wbar}
\end{align}
由于\eqref{align:ch15-smulian-approx}对所有 $n$ 成立,
\begin{align}
\text{若 }n\leqslant N(m-1),\quad|(T-J(x_{m}))[\psi_{n}]|<1/m.\label{align:ch15-smulian-weak}
\end{align}
由于 $B$ 是序列紧的, $\{x_{n}\}$ 存在弱收敛于 $x\in B$ 的子序列 $\{x_{n_{k}}\}$. Mazur 定理告诉我们序列 $\{x_{n_{k}}\}$ 中的项的凸组合的序列强收敛于 $x$. 该凸组合序列在 $J$ 下的象在 $X^{**}$ 中强收敛于 $J(x)$. 因此 $J(x)$ 属于 $W$. 但 $T$ 也属于 $W$. 因此 $T-J(x)$ 属于 $W$. 我们宣称 $T=J(x)$. 根据\eqref{align:ch15-smulian-Wbar}, 为证明该断言必要与充分, 要证明:
\begin{align}
\text{对所有 }n,\quad(T-J(x))[\psi_{n}]=0.\label{align:ch15-smulian-zero}
\end{align}
固定自然数 $n$. 观察到对每个指标 $k$,
\begin{align*}
(T-J(x))[\psi_{n}]=(T-J(x_{n_{k}}))[\psi_{n}]+(J(x_{n_{k}})-J(x))[\psi_{n}].
\end{align*}
我们从\eqref{align:ch15-smulian-weak}推出, 若 $N(n_{k}-1)>n$, 则 $|(T-J(x_{n_{k}}))[\psi_{n}]|<1/n_{k}$. 另一方面,
\begin{align*}
\text{对所有 }k,\quad(J(x_{n_{k}})-J(x))[\psi_{n}]=\psi_{n}(x_{n_{k}}-x)
\end{align*}
且 $\{x_{n_{k}}\}$ 弱收敛于 $x$. 因此,
\begin{align*}
(T-J(x))[\psi_{n}]=\lim_{k\to\infty}(T-J(x_{n_{k}}))[\psi_{n}]+\lim_{k\to\infty}(J(x_{n_{k}})-J(x))[\psi_{n}]=0.
\end{align*}
\end{proof}

我们把 Kakutani 定理与 Eberlein-\v{S}mulian 定理概括为以下陈述.

\begin{theorem}[弱紧性的刻画]\label{theorem:ch15-weak-compactness}
令 $B$ 为 Banach 空间 $X$ 的闭单位球. 则以下三个断言等价:
\begin{enumerate}[(i)]
\item $X$ 是自反的;
\item $B$ 是弱紧的;
\item $B$ 是弱序列紧的.
\end{enumerate}
\end{theorem}

\begin{exercise}\label{exercise:ch15-13}
在不可度量化的一般拓扑空间中, 序列可能收敛于不止一个点. 证明对赋范线性空间上的 $W$ 弱拓扑, 其中 $W$ 是 $X^{*}$ 的分离 $X$ 中的点的子空间, 这不可能发生.
\end{exercise}

\begin{exercise}\label{exercise:ch15-14}
证明 $L^{\infty}[0,1]$ 中存在没有弱收敛子序列的有界序列. 证明 $C[a,b]$ 的闭单位球不是弱紧的.
\end{exercise}

\begin{exercise}\label{exercise:ch15-15}
令 $K$ 为具有无穷多个点的紧度量空间. 证明 $C(K)$ 中存在没有弱收敛子序列的有界序列(见习题\cref{exercise:ch15-11}), 但 $C(K)$ 上的每个连续线性泛函的有界序列有逐点收敛到 $C(K)$ 上的连续线性泛函的子序列.
\end{exercise}

\end{document}